\scanpage{285}
\mspage{25}

peut être identifié au monoïde libre \inserted{$\mathrm{Ty}$} à deux générateurs $\Delta^+$ et $\Delta^-$, dont les éléments \deleted{s'obtiennent en faisant} sont les suites finies quelconques (evt vides) d'objets alternativement \unsure{égales|égaux} à $\Delta^+$ et à $\Delta^-$. \deleted{\unsure{Un}}

Si le premier élément est $\Delta^+$, on a un \deleted{\unsure{a}} type de chemin à di\uline{rec}tion initiale directe, rétrograde dans le cas contraire. La décomposition de $I_\tau$ \inserted{ou de $\mathrm{Fl}(I_\tau)$} en intervalles critiques correspond à l'écriture réduite d'un mot comme
\[ {\Delta^+}^{\ell_1}(\Delta^-)^{\ell_2}\ldots(\Delta^{\unsure{?}})^{\ell_{\unsure{r|v}}}
\quad\left(\begin{aligned}\Delta^{?}&=\Delta^+&&\text{si }r\text{ impair}\\\Delta^{?}&=\Delta^-&&\text{si }r\text{ pair}\end{aligned}\right) \]
ou comme
\[ (\Delta^-)^{\ell_1}(\Delta^+)^{\ell_2}\ldots(\Delta^{?})^{\ell_r}
\quad\left(\begin{aligned}\Delta^{?}&=\Delta^-&&\text{si }r\text{ impair}\\
\deleted{\unsure{\delta}}\Delta^{?}&=\Delta^+&&\text{si }r\text{ pair}\end{aligned}\right). \]
Les $\ell_i$ sont les longueurs des intervalles critiques successifs (ou des chaînes structurales successives), leur nombre est la longueur réduite de l'intervalle.

On va considérer $\mathrm{Ty}$ comme l'ens des objets d'une catégorie $\underline{\mathrm{Ty}}$
\[ \left\{\begin{aligned}
\mathrm{Ob}\,\underline{\mathrm{Ty}}&=\mathrm{Ty}=\begin{aligned}[t]
&\text{\deleted{ensemble des} monoïde libre à deux}\\
&\text{générateurs }\inserted{\Delta^+,\Delta^-}\ \text{\unsure{précisés|précédents}.}
\end{aligned}\\
\text{Si }\tau,\tau'&\in\mathrm{Ty},\quad \mathrm{Hom}(\tau,\tau')=\mathrm{Hom}_{\unsure{\mathrm{Cat}}}(I_\tau,I_{\tau'}).
\end{aligned}\right. \]
% slanted
\margin{Définition un peu \unsure{cavalière}, il faut ne prendre que les $\varphi$ de $I_\tau$ dans $I_{\tau'}$ compatibles avec les \deleted{\unsure{objets}} sources et buts --- cf \unsure{plus bas|p. suiv.}.}

\unsure{On notera qu'ici|On notera qu'on a|On peut même préciser}, \deleted{\unsure{non}} seulement un ensemble $\mathrm{Hom}_{\mathrm{Cat}}(I_\tau,I_{\tau'})$, mais même une \uline{catégorie} $\underline{\mathrm{Hom}}(I_\tau,I_{\tau'})$, ce qui revient ici à dire que $\mathrm{Hom}(I_\tau,I_{\tau'})$ (\unsure{i.e.} $\mathrm{Hom}(\tau,\tau')$) peut être vu comme un ens. \uline{ordonné}. Mais pour l'instant nous allons ignorer cette relation d'ordre.
% archive sheet number: 284
